\chapter{Prediction and evolutionary analysis of gene function}

%%%%%%%%%%%%%%%%%%%%%%
%%%%%%%%%%%%%%%%%%%%%%
\section{Introduction}
%%%%%%%%%%%%%%%%%%%%%%
%%%%%%%%%%%%%%%%%%%%%%


	Characterizing the functional behavior of individual proteins in a variety of different contexts is an important first step towards understanding life at the molecular level.  Endeavors such as understanding biological pathways, investigating disease, and developing drugs to cure those diseases depend on being able to describe the actions of individual proteins, both in terms of their physiochemical molecular function, involvement in biological processes, and the sub-cellular location at which these actions are carried out.
	
		give some example of how this individualized approach has made progress.  find a disease.  maybe speak about lacz
	
	
	In spite of the fact that there are increasingly more tools for the interrogation of protein behavior, there are still a large number of functionally uncharacterized proteins, with the gap between known sequences and experimentally characterized ones growing at an exponential pace (\Cref{growth}). 
	
	Furthermore, while most experimentally annotated proteins come from model organisms, with the exception of yeast, less than half of the genome in any model organisms has yet to be assigned experimentally characterized GO functions (\Cref{per_organism}). 
	
	An even more disparaging point is that the numbers in \Cref{per_organism}, is that they do not fully represent the vagueness of many annotations: \textcolor{red}{$41\%$} of experimentally characterized proteins in the \textcolor{red}{January 2011 version of SwissProt} have �protein binding� as their most specific molecular function term. Given the multitude of ways a protein's function can be characterized, having a single annotation does not mean that protein has been fully annotated.
	

	
	
	
	Although a growing number of tools for the task of interrogating a protein's function are at the disposal of experimental biologists, characterizing the exponentially expanding set of known sequences requires the application of \textit{in silico} techniques. 
	
	
Current techniques for the evaluation of function prediction are reviewed and several drawbacks are addressed through the introduction of new information content-based metrics that attempt to provide a comprehensive assessment of the performance of a given prediction method. Finally, the influence of speciation and gene duplication events on the evolution of protein function is addressed by comparing the functional similarity of differing classes of homologous sequences.  





Characterizing the function of individual proteins is important 


Experimental studies of protein function have been steadily accumulating over the years. Currently, there are about $50,000$ proteins with at least one experimentally annotated GO term from MFO or BPO in the GO database and Swiss-Prot \citep{Bairoch2005} combined. However, owing to the numerous sequencing projects \citep{Liolios2008}, the gap between annotated and non-annotated proteins has exceeded two orders of magnitude, and will only get wider. In addition, annotated proteins may not be sufficiently covered, both with respect to the resolution of functional annotations and the inclusion of other yet unknown functions. Thus, it is important to develop algorithms capable of accurately predicting protein function.





\section{Defining function}



Give a brief overview of the various definitions of function.  FOcus on the selected effect and causal role definitions of function.  

Categorize the use of function in the context of the gene ontology.

Distance yourself from any teleological explanations.  


Protein function is a multi-scale concept that reflects ``everything that happens to or through a protein"\citep{Rost2003}, and is typically considered from the biochemical, biological, and phenotypic perspectives \citep{Bartlett2003}. From the biochemical, or molecular, perspective a protein may be a kinase, whereas in terms of its biological function this kinase can be involved in numerous processes, such as cell cycle regulation or cell-cell signaling. Two proteins with the same molecular function may be involved in drastically different biological processes, and conversely, the set of proteins associated with a particular biological process will generally be drawn from a wide range of molecular functions. From the phenotypic viewpoint, a protein is generally associated with variation in observable physical or behavioral traits. For example, kinase variants or mutants may be responsible for disease. Adding another level of complexity to the study of function is the fact that a protein's molecular function, while generally considered to be a static notion, is modulated by a particular cellular context, e.g. the presence of other molecules, or properties of the physical environment, e.g. temperature \citep{Mohan2009}.

Several classification systems have been proposed to standardize functional annotation and to facilitate computation. With few exceptions, these classification systems usually take on the structure of hierarchical ontologies. Enzyme Commission (EC) numbers \citep{NC_IUBMB_1992} and the MIPS functional catalogue \citep{Ruepp2004} are two well-accepted schemes; however, the most commonly used functional classification is the Gene Ontology (GO). GO provides three hierarchical classifications as directed acyclic graphs: Molecular Function Ontology (MFO), Biological Process Ontology (BPO), and Cellular Component Ontology \citep{Ashburner2000}. With respect to defining a particular gene's phenotype, the existing classifications are species-specific (e.g. the Human Phenotype Ontology \citep{Robinson2010}) and predominantly constructed to address human disease. The Unified Medical Language System \citep{Bodenreider2004}, for example, incorporates a number of vocabularies together with semantic relationships between terms, mainly for the purpose of defining associations between genes and medical disorders.

This dissertation addresses protein function at the molecular function and biological process levels, as defined by the GO consortium, and functional inference from protein sequence alone. We analyze the power of functional transfer of GO terms using sequence alignments and develop a method for probabilistic inference of GO terms using supervised learning. Our algorithm, Functional ANNotator (FANN), employs multi-output artificial neural networks. We show that in the GO annotation task, FANN-GO outperforms standard sequence alignment methods and GOtcha.

\section{why predicting function is necessary?}

I mean, is it really?  Are we just waiting compute time on the computing cluster? Are the predictions produced, if correct, even useful?

\section{State of the art methods}

Historically, sequence-based inference was the first strategy used to predict protein function, even if most studies at the time avoided explicitly relating homology and function \citep{Doolittle1986}. Global and local sequence alignments were used to query sequence databases for similarities with a target protein. With the accumulation of experimentally determined protein functions, the most similar annotated sequences have traditionally been used to infer function \citep{Rost2003}. More advanced methods exploited predicted physicochemical properties \citep{Jensen2002,Jensen2003}, evolutionary relationships\citep{Engelhardt2005,Pellegrini1999,Bandyopadhyay2006b}, or the structure of functional ontologies in order to achieve different confidence levels at different ontological terms \citep{Barutcuoglu2006,Hawkins2006,Martin2004}. Microarrays, protein-protein interaction networks, protein structures, and predicted ligands have also been exploited \citep{Deng2003,Hermann2007,Lee2004,Troyanskaya2003,Letovsky2003,Nabieva2005,Pal2005,Pazos2004,Song2007,Vazquez2003}. However, most of these methods are limited to a few organisms where such data are available. One way or another, sequence alignment-based inference is the cornerstone of functional inference and is the focal point of this study.

Sequence alignment-based transfer of function has been thoroughly studied in the last decade, predominantly for enzymes \citep{Addou2009,Devos2000,Rost2003,Tian2003,Todd2001,Wilson2000}. The results of these studies indicate that at least $60\%$ sequence identity, and more likely closer to $80\%$, is required for the accurate transfer of the third level of EC classification. More sophisticated approaches were proposed as well: the GOtcha method was developed in order to take sequence alignment scores between a query protein and a functionally annotated database and overlay them on the functional ontology, cumulatively propagating such scores \citep{Martin2004}. PFP refined this technique by incorporating PSI-BLAST alignments at very low significance levels and conditional probabilities that a protein is associated with pairs of functional terms \citep{Hawkins2006} . Other methods such as ProtFun \citep{Jensen2003}, ConFunc \citep{Wass2008}, GOsling \citep{Jones2008a}, and GOstruct \citep{Sokolov2010} were developed for high-throughput prediction tasks. Finally, phylogenetic methods attempt to exploit particular evolutionary relationships within a gene family \citep{Brown2006,Eisen1998}. Methods such as SIFTER \citep{Engelhardt2005} or ortholog identification methods \citep{Remm2001} belong in this category. Several recent reviews provide good perspectives on protein function prediction at all scales \citep{Dalkilic2008,Friedberg2006,Kann2007,Laskowski2008,Lee2007,Punta2008,Rentzsch2009,Rost2003}.

\subsection{Community wide assessment of function prediction}

\section{Analysis of function in evolutionary context}

The previous chapters focused largely on the inference of function.  While these endeavors are important in their own right, It is desirable that the data produced by these efforts will at some point serve as more than an encyclopedic cataloging of the known gene/protein universe.


In \Cref{ch:oc} 

\section{Conclusion}



list contributions

